\chapter{Clase 7}

\section{Funciones de varias variables}

En el estudio de múltiples variables aleatorias, es fundamental poder calcular el valor esperado de una función $h(X,Y)$ que depende de dos variables aleatorias $X$ y $Y$. Este valor esperado se obtiene ponderando la función por la probabilidad conjunta de las variables (Devore, Capítulo 5, Sección 5.2).

Para variables aleatorias conjuntamente discretas con función de masa de probabilidad conjunta $p(x,y)$:
$$ E[h(X,Y)] = \sum_{x} \sum_{y} h(x,y) p(x,y) $$

Para variables aleatorias conjuntamente continuas con función de densidad de probabilidad conjunta $f(x,y)$:
$$ E[h(X,Y)] = \int_{-\infty}^{\infty} \int_{-\infty}^{\infty} h(x,y) f(x,y) \, dx \, dy $$

Una propiedad fundamental derivada de esta definición es la linealidad del valor esperado. Si la función es una combinación lineal $h(X,Y) = aX + bY$, donde $a$ y $b$ son constantes, entonces:
$$ E(aX + bY) = aE(X) + bE(Y) = a\mu_X + b\mu_Y $$

\section{Distribuciones de probabilidad conjunta}

El comportamiento simultáneo de dos variables aleatorias se describe mediante su distribución conjunta. 

Para el caso discreto, la función de masa de probabilidad conjunta $p(x,y)$ especifica la probabilidad de que $X=x$ y $Y=y$ simultáneamente. Debe satisfacer rigurosamente (Devore, Capítulo 5, Sección 5.1):
$$ p(x,y) \ge 0 \quad \text{para todo } (x,y) $$
$$ \sum_{x} \sum_{y} p(x,y) = 1 $$

Para el caso continuo, la función de densidad de probabilidad conjunta $f(x,y)$ define la probabilidad de que $(X,Y)$ caiga en una región bidimensional $A$. Sus condiciones analíticas son (Devore, Sección 5.1):
$$ f(x,y) \ge 0 \quad \text{para todo } (x,y) $$
$$ \int_{-\infty}^{\infty} \int_{-\infty}^{\infty} f(x,y) \, dx \, dy = 1 $$
$$ P((X,Y) \in A) = \iint_{A} f(x,y) \, dx \, dy $$

\section{Distribuciones marginales}

Las distribuciones marginales describen el comportamiento probabilístico de una sola variable (por ejemplo, $X$), ignorando o promediando los resultados de la otra variable ($Y$). Se obtienen directamente de la distribución conjunta (Devore, Capítulo 5, Sección 5.1).

Para variables discretas, la función de masa de probabilidad marginal de $X$ se calcula sumando la probabilidad conjunta sobre todos los posibles valores de $Y$:
$$ p_X(x) = \sum_{y} p(x,y) $$

Para variables continuas, la función de densidad de probabilidad marginal de $X$ se obtiene integrando la densidad conjunta sobre todo el rango de $Y$:
$$ f_X(x) = \int_{-\infty}^{\infty} f(x,y) \, dy $$

De manera simétrica, las distribuciones marginales para $Y$ se definen como $p_Y(y) = \sum_{x} p(x,y)$ en el caso discreto, y $f_Y(y) = \int_{-\infty}^{\infty} f(x,y) \, dx$ en el caso continuo.

\section{Combinación de variables aleatorias}

Al combinar variables aleatorias, es crucial cuantificar no solo sus medias individuales, sino también cómo varían conjuntamente. La covarianza mide la relación lineal entre $X$ y $Y$, y se define como el valor esperado del producto de sus desviaciones respecto a sus respectivas medias (Devore, Capítulo 5, Sección 5.2):
$$ Cov(X,Y) = E[(X-\mu_X)(Y-\mu_Y)] = E(XY) - \mu_X\mu_Y $$

El coeficiente de correlación $\rho$ estandariza la covarianza para acotarla entre $-1$ y $1$:
$$ \rho_{X,Y} = \frac{Cov(X,Y)}{\sigma_X \sigma_Y} $$

Para una combinación lineal de variables aleatorias $U = aX + bY$, la varianza de $U$ depende de las varianzas individuales y de su covarianza. La expresión general es (Devore, Capítulo 5, Sección 5.5):
$$ V(aX + bY) = a^2V(X) + b^2V(Y) + 2abCov(X,Y) $$

Si las variables aleatorias $X$ y $Y$ son estadísticamente independientes, su covarianza es nula ($Cov(X,Y) = 0$). En este escenario, la varianza de la combinación lineal se simplifica a la suma ponderada de las varianzas individuales:
$$ V(aX + bY) = a^2V(X) + b^2V(Y) $$


\subsection{Ejemplo 1: Abundancia Mineralógica en una Sección Delgada de Meteorito}

Principios de petrología y análisis de secciones delgadas expuestos en \textit{Planetary Geology} (Capítulo sobre \textit{Mineralogy and Petrology of Meteorites}). 

Un laboratorio de geoquímica isotópica analiza una sección delgada de una condrita ordinaria. El microscopio electrónico mide las fracciones normalizadas de dos minerales de silicato fundamentales: Olivino (variable aleatoria $X$) y Piroxeno (variable aleatoria $Y$). Debido a la presencia de matriz vítrea y otros minerales accesorios, la suma de las fracciones de Olivino y Piroxeno no puede exceder la unidad. 

La función de densidad de probabilidad conjunta que modela la ocurrencia de estas abundancias en esta clase específica de meteoritos está dada por:
$$
f(x,y) = 
\begin{cases} 
c(x + y) & \text{para } x > 0, y > 0, x + y < 1 \\
0 & \text{en otro caso}
\end{cases}
$$

A continuación, resolvemos una secuencia de requerimientos para caracterizar completamente el sistema mineralógico:

\textbf{a) Determinación de la constante de normalización $c$}
Para que $f(x,y)$ sea una FDP legítima, la integral doble sobre toda la región de soporte debe ser igual a 1 (Devore, Capítulo 5, Sección 5.1). La región es un triángulo en el primer cuadrante acotado por $x+y=1$:
$$ \int_{0}^{1} \int_{0}^{1-x} c(x + y) \, dy \, dx = 1 $$
Resolvemos la integral interna respecto a $y$:
$$ \int_{0}^{1-x} (x + y) \, dy = \left[ xy + \frac{y^2}{2} \right]_{0}^{1-x} = x(1-x) + \frac{(1-x)^2}{2} $$
Simplificando la expresión algebraica:
$$ x - x^2 + \frac{1 - 2x + x^2}{2} = \frac{2x - 2x^2 + 1 - 2x + x^2}{2} = \frac{1 - x^2}{2} $$
Ahora integramos respecto a $x$:
$$ c \int_{0}^{1} \frac{1 - x^2}{2} \, dx = \frac{c}{2} \left[ x - \frac{x^3}{3} \right]_{0}^{1} = \frac{c}{2} \left( 1 - \frac{1}{3} \right) = \frac{c}{2} \left( \frac{2}{3} \right) = \frac{c}{3} $$
Igualando a 1, obtenemos $c = 3$. La FDP conjunta es $f(x,y) = 3(x+y)$ en la región definida.

\textbf{b) Obtención de la distribución marginal de $X$}
La distribución marginal de $X$ describe la probabilidad de encontrar una cierta abundancia de Olivino, independientemente de la cantidad de Piroxeno. Se obtiene integrando la conjunta sobre todos los valores posibles de $Y$ (Devore, Sección 5.1):
$$ f_X(x) = \int_{0}^{1-x} 3(x + y) \, dy $$
Utilizando el resultado de la integral interna calculado en el paso anterior:
$$ f_X(x) = 3 \left( \frac{1 - x^2}{2} \right) = \frac{3}{2}(1 - x^2) \quad \text{para } 0 < x < 1 $$

\textbf{c) Valor esperado de la fracción total de silicatos}
El índice de saturación de silicatos se define como la suma de las fracciones, $Z = X + Y$. Calculamos su valor esperado utilizando la definición para funciones de variables conjuntas (Devore, Sección 5.2):
$$ E(X + Y) = \int_{0}^{1} \int_{0}^{1-x} (x + y) \cdot 3(x + y) \, dy \, dx = 3 \int_{0}^{1} \int_{0}^{1-x} (x + y)^2 \, dy \, dx $$
Resolvemos la integral interna haciendo el cambio mental $u = x+y$:
$$ \int_{0}^{1-x} (x + y)^2 \, dy = \left[ \frac{(x+y)^3}{3} \right]_{y=0}^{y=1-x} = \frac{1^3}{3} - \frac{x^3}{3} = \frac{1 - x^3}{3} $$
Sustituimos en la integral externa:
$$ E(X + Y) = 3 \int_{0}^{1} \frac{1 - x^3}{3} \, dx = \int_{0}^{1} (1 - x^3) \, dx = \left[ x - \frac{x^4}{4} \right]_{0}^{1} = 1 - \frac{1}{4} = \frac{3}{4} $$
El valor esperado de la fracción total de silicatos es 0.75.

\textbf{d) Covarianza y Correlación entre Olivino y Piroxeno}
Para calcular la covarianza, necesitamos $E(X)$, $E(Y)$ y $E(XY)$. Por simetría de la función $f(x,y)$ respecto a $X$ y $Y$, sabemos que $E(X) = E(Y)$.
$$ E(X) = \int_{0}^{1} x f_X(x) \, dx = \int_{0}^{1} x \frac{3}{2}(1 - x^2) \, dx = \frac{3}{2} \int_{0}^{1} (x - x^3) \, dx = \frac{3}{2} \left[ \frac{x^2}{2} - \frac{x^4}{4} \right]_{0}^{1} = \frac{3}{2} \left( \frac{1}{4} \right) = \frac{3}{8} $$
Calculamos $E(X^2)$ para hallar la varianza de $X$:
$$ E(X^2) = \int_{0}^{1} x^2 \frac{3}{2}(1 - x^2) \, dx = \frac{3}{2} \left[ \frac{x^3}{3} - \frac{x^5}{5} \right]_{0}^{1} = \frac{3}{2} \left( \frac{2}{15} \right) = \frac{1}{5} $$
$$ V(X) = E(X^2) - [E(X)]^2 = \frac{1}{5} - \left(\frac{3}{8}\right)^2 = \frac{1}{5} - \frac{9}{64} = \frac{64 - 45}{320} = \frac{19}{320} $$
Por simetría, $V(Y) = 19/320$. Ahora calculamos $E(XY)$:
$$ E(XY) = \int_{0}^{1} \int_{0}^{1-x} xy \cdot 3(x + y) \, dy \, dx = 3 \int_{0}^{1} \int_{0}^{1-x} (x^2y + xy^2) \, dy \, dx $$
Integral interna:
$$ \left[ \frac{x^2y^2}{2} + \frac{xy^3}{3} \right]_{0}^{1-x} = \frac{x^2(1-x)^2}{2} + \frac{x(1-x)^3}{3} = \frac{x(1-x)^2}{6} [3x + 2(1-x)] = \frac{x(1-x)^2(2+x)}{6} $$
Expandiendo el polinomio: $\frac{1}{6}(x - 2x^2 + x^3)(2+x) = \frac{1}{6}(2x - 3x^2 + x^4)$.
Integral externa:
$$ E(XY) = 3 \int_{0}^{1} \frac{1}{6}(2x - 3x^2 + x^4) \, dx = \frac{1}{2} \left[ x^2 - x^3 + \frac{x^5}{5} \right]_{0}^{1} = \frac{1}{2} \left( 1 - 1 + \frac{1}{5} \right) = \frac{1}{10} $$
Finalmente, la covarianza y el coeficiente de correlación (Devore, Sección 5.2):
$$ Cov(X,Y) = E(XY) - E(X)E(Y) = \frac{1}{10} - \left(\frac{3}{8}\right)\left(\frac{3}{8}\right) = \frac{32}{320} - \frac{45}{320} = -\frac{13}{320} $$
$$ \rho_{X,Y} = \frac{Cov(X,Y)}{\sqrt{V(X)V(Y)}} = \frac{-13/320}{19/320} = -\frac{13}{19} \approx -0.684 $$
La correlación negativa fuerte es físicamente consistente: al estar limitadas por $X+Y < 1$, una alta abundancia de Olivino restringe matemáticamente la cantidad máxima de Piroxeno que puede coexistir en la sección delgada.

\textbf{e) Varianza de un índice petrográfico combinado}
El laboratorio define un Índice de Diferenciación Petrográfica como $W = 2X + 3Y$. Utilizando la regla de la varianza para combinaciones lineales (Devore, Sección 5.5):
$$ V(W) = 2^2 V(X) + 3^2 V(Y) + 2(2)(3) Cov(X,Y) $$
$$ V(W) = 4\left(\frac{19}{320}\right) + 9\left(\frac{19}{320}\right) + 12\left(-\frac{13}{320}\right) = \frac{76 + 171 - 156}{320} = \frac{91}{320} \approx 0.284 $$
Este resultado cuantifica rigurosamente cómo la incertidumbre en las mediciones minerales individuales y su correlación negativa se propagan en el índice compuesto, permitiendo a los geólogos planetarios establecer márgenes de confianza para la clasificación del meteorito.

\subsection{Ejemplo 2: Detección Multi-banda de Tránsitos Exoplanetarios}

En la astronomía observacional, la confirmación de exoplanetas requiere cruzar datos de diferentes instrumentos o bandas espectrales para descartar falsos positivos. Inspirados en los métodos de fotometría y validación estadística de \textit{Astronomical Measurement} (Capítulo sobre \textit{Transit Photometry and False Positive Probabilities}), modelaremos un escenario de detección discreta.

Un telescopio espacial analiza un campo estelar. Para una estrella objetivo específica, definimos las variables aleatorias discretas:
\begin{itemize}
    \item $X$: Número de señales candidatas a tránsito detectadas en la banda visual.
    \item $Y$: Número de señales candidatas a tránsito detectadas en la banda infrarroja.
\end{itemize}

Basado en el rendimiento histórico del pipeline de reducción de datos, la función de masa de probabilidad conjunta $p(x,y)$ para $X \in \{0, 1, 2\}$ y $Y \in \{0, 1\}$ se resume en la siguiente tabla:

$$
\begin{array}{c|ccc}
 & X=0 & X=1 & X=2 \\
\hline
Y=0 & 0.10 & 0.15 & 0.05 \\
Y=1 & 0.20 & 0.25 & 0.25
\end{array}
$$

A continuación, desarrollamos los pasos para evaluar la estadística de detección del sistema:

\textbf{a) Distribuciones marginales}
Las distribuciones marginales describen el comportamiento de cada banda por separado, sumando las probabilidades conjuntas (Devore, Sección 5.1):
Para $X$:
$$ p_X(0) = 0.10 + 0.20 = 0.30 $$
$$ p_X(1) = 0.15 + 0.25 = 0.40 $$
$$ p_X(2) = 0.05 + 0.25 = 0.30 $$
Para $Y$:
$$ p_Y(0) = 0.10 + 0.15 + 0.05 = 0.30 $$
$$ p_Y(1) = 0.20 + 0.25 + 0.25 = 0.70 $$

\textbf{b) Valores esperados y varianzas individuales}
Calculamos los momentos de primer y segundo orden para cada variable (Devore, Capítulo 3, Sección 3.3):
$$ E(X) = 0(0.30) + 1(0.40) + 2(0.30) = 1.0 $$
$$ E(X^2) = 0^2(0.30) + 1^2(0.40) + 2^2(0.30) = 0.40 + 1.20 = 1.60 $$
$$ V(X) = E(X^2) - [E(X)]^2 = 1.60 - 1.0^2 = 0.60 $$

$$ E(Y) = 0(0.30) + 1(0.70) = 0.70 $$
$$ E(Y^2) = 0^2(0.30) + 1^2(0.70) = 0.70 $$
$$ V(Y) = E(Y^2) - [E(Y)]^2 = 0.70 - 0.70^2 = 0.70 - 0.49 = 0.21 $$

\textbf{c) Covarianza y Correlación}
Para evaluar si las detecciones en ambas bandas son independientes o están correlacionadas, calculamos $E(XY)$ sumando sobre todos los pares $(x,y)$ donde la probabilidad es distinta de cero:
$$ E(XY) = \sum_{x} \sum_{y} xy \cdot p(x,y) $$
Los únicos términos que no se anulan (donde $x \neq 0$ y $y \neq 0$) son:
$$ E(XY) = (1)(1)(0.25) + (2)(1)(0.25) = 0.25 + 0.50 = 0.75 $$
La covarianza es:
$$ Cov(X,Y) = E(XY) - E(X)E(Y) = 0.75 - (1.0)(0.70) = 0.05 $$
El coeficiente de correlación es:
$$ \rho_{X,Y} = \frac{Cov(X,Y)}{\sqrt{V(X)V(Y)}} = \frac{0.05}{\sqrt{0.60 \times 0.21}} = \frac{0.05}{\sqrt{0.126}} \approx \frac{0.05}{0.355} \approx 0.141 $$
La correlación positiva débil sugiere que, aunque las bandas son físicamente distintas, existen factores astrofísicos subyacentes (como la actividad estelar o el tamaño del planeta) que tienden a aumentar la tasa de detección en ambas bandas simultáneamente.

\textbf{d) Varianza de un discriminante de validación}
Para separar los planetas reales de los binarios eclipsantes de fondo, el pipeline calcula un Puntaje Discriminante $Z = 3X - 2Y$. Las ponderaciones positivas para la banda visual y negativas para la infrarroja ayudan a filtrar ciertos tipos de falsos positivos. Necesitamos la varianza de este puntaje para calibrar los umbrales de alerta.

Aplicando la fórmula de la varianza para una combinación lineal (Devore, Sección 5.5):
$$ V(Z) = V(3X - 2Y) = 3^2 V(X) + (-2)^2 V(Y) + 2(3)(-2) Cov(X,Y) $$
$$ V(Z) = 9(0.60) + 4(0.21) - 12(0.05) $$
$$ V(Z) = 5.40 + 0.84 - 0.60 = 5.64 $$

La varianza del puntaje discriminante es 5.64. Conocer este valor exacto permite a los científicos de datos establecer los intervalos de confianza para el puntaje $Z$. Si la varianza hubiera sido calculada asumiendo erróneamente que $X$ e $Y$ eran independientes (ignorando el término de covarianza), el resultado habría sido $6.24$, lo que llevaría a subestimar la precisión del discriminante y a un diseño ineficiente de los filtros de validación automática. Este ejercicio demuestra el poder de las distribuciones conjuntas para modelar sistemas de detección complejos en la astronomía moderna.
